\section{Worked example}
\label{sec:worked-example}

\begin{keynote}
\noindent\textbf{This example is illustrative. The site does not exist, and the
result is not a finding about any real place.}

The calculation was performed \textbf{by hand} from the formulas and
configuration values specified in this paper, and every number below is also
asserted, digit for digit, by the engine's golden-scenario suite --- the hand
derivation is the expected value, never the engine's own output. Its purpose
is to make the specification concrete enough to check: a reader who disagrees
with any number below can locate the equation that produced it and recompute
it.

One quantity still rests on an assumption: the grid emission factor of
\Cref{sec:ghg}, which no shipped configuration supplies.
\Cref{sec:we-assumptions} states how much the conclusion depends on it.
\end{keynote}

\subsection{The illustrative site and intervention}

A neighbourhood-scale street corridor in a dense, largely sealed inner-city
area in a temperate climate. The proposed intervention is street tree planting
(typology 01). \Cref{tab:we-inputs} gives the inputs.

\begin{table}[htbp]
\centering
\caption[Worked example inputs]{Inputs to the illustrative assessment. Fields
left unsupplied are shown explicitly, because their absence changes both the
result and its confidence.}
\label{tab:we-inputs}
\small
\begin{tabularx}{\textwidth}{@{}l>{\raggedright\arraybackslash}Xl@{}}
\toprule
\textbf{Group} & \textbf{Field} & \textbf{Value} \\
\midrule
Project    & Assessment scale                     & Neighbourhood \\
\addlinespace
Site       & Site area                            & \qty{6000}{\metre\squared} \\
           & Existing tree canopy                 & \SI{8}{\percent} \\
           & Existing green cover                 & \SI{12}{\percent} \\
           & Impervious surface                   & \SI{85}{\percent} \\
           & Soil availability                    & Limited \\
           & Irrigation availability              & Occasional \\
\addlinespace
Climate    & Climate zone                         & Temperate \\
           & Land surface temperature anomaly     & \qty{+4.2}{\degreeCelsius} \\
           & Solar exposure                       & High \\
\addlinespace
Vulnerability & Population density                & High \\
           & Vulnerable population presence       & High \\
           & Access to cooled indoor space        & Low \\
\addlinespace
Intervention & Typology                           & Street tree planting (01) \\
           & New canopy area at maturity          & \qty{1200}{\metre\squared} \\
\addlinespace
Energy     & Nearby building cooling demand relevant & Yes \\
           & Annual cooling energy demand         & \qty{450000}{\kilo\watt\hour} \\
           & Energy price                         & \emph{not supplied} \\
\addlinespace
Cost       & Capital cost                         & \emph{not supplied} \\
\bottomrule
\end{tabularx}
\end{table}

\subsection{Layer 1 --- Baseline assessment}

An LST anomaly is available, so the data-rich path of
\Cref{eq:heat-exposure-rich} applies. Normalising by \Cref{eq:lst},
$S_{\mathrm{LST}} = \clamp{10 \times 4.2} = 42.0$. The impervious percentage
enters directly as $S_{\mathrm{imp}} = 85$; solar exposure \emph{high}
normalises to $S_{\mathrm{solar}} = 75$ by \Cref{tab:levels}; and the
vegetation deficit is $S_{\mathrm{vegdef}} = \clamp{100 - 12} = 88$ by
\Cref{eq:vegdef}. Then

\begin{align*}
S_{\mathrm{heat}}
&= 0.40(42.0) + 0.25(85) + 0.20(75) + 0.15(88) \\
&= 16.80 + 21.25 + 15.00 + 13.20 \;=\; \mathbf{66.25}.
\end{align*}

For vulnerability, population density and vulnerable-group presence are both
\emph{high} and normalise to 75. Access to cooled indoor space is \emph{low},
which under the inverted mapping of \Cref{sec:inverted} yields a cooling access
\emph{deficit} of 100. By \Cref{eq:vulnerability},

\[
S_{\mathrm{vuln}} = 0.40(75) + 0.40(75) + 0.20(100)
                  = 30.00 + 30.00 + 20.00 = \mathbf{80.00}.
\]

The Heat Priority Index follows from \Cref{eq:hpi}:

\[
\mathrm{HPI} = 0.60(66.25) + 0.40(80.00) = 39.75 + 32.00 = \mathbf{71.75},
\]

which falls in the \emph{High} band of \Cref{tab:hpi-bands}. Note that
vulnerability (80.0) exceeds heat exposure (66.25) here, so the equity-forward
weighting of \Cref{sec:effective-weights} raises this site's final score
relative to a heat-only instrument. That is the disclosed value choice
operating as intended, and a reviewer who rejects the choice should read this
result as correspondingly inflated.

\subsection{Layer 2 --- NbS performance}

Each of the four adjustment conditions of \Cref{tab:derivation} is derived from
inputs already supplied.

\emph{Canopy.} Existing tree canopy is \SI{8}{\percent} of
\qty{6000}{\metre\squared}, or \qty{480}{\metre\squared}. Adding
\qty{1200}{\metre\squared} of new canopy at maturity gives
\qty{1680}{\metre\squared}, which is \SI{28.0}{\percent} of the site --- inside
the \emph{good} band ($25 \le c < 40$), so $f_{\mathrm{canopy}} = 1.0$. The site
falls short of the \SI{40}{\percent} inflection identified by
\textcite{ziter2019}, which is why it does not reach \emph{excellent}.

\emph{Soil--water.} Soil availability \emph{limited} maps to \emph{moderate};
irrigation availability \emph{occasional} also maps to \emph{moderate}. The
more limiting of the two is \emph{moderate}, so
$f_{\mathrm{soilwater}} = 0.8$.

\emph{Scale.} Assessment scale \emph{neighbourhood} maps to \emph{good}, so
$f_{\mathrm{scale}} = 1.0$.

\emph{Climate.} Street tree planting is a \emph{green}-family typology in a
\emph{temperate} zone, which \Cref{tab:climate-matrix} rates \emph{good}, so
$f_{\mathrm{climate}} = 1.0$.

By \Cref{eq:adjustment},

\[
A = 0.40(1.0) + 0.25(0.8) + 0.20(1.0) + 0.15(1.0)
  = 0.400 + 0.200 + 0.200 + 0.150 = \mathbf{0.95}.
\]

Street tree planting has base cooling score $B = 75$ and envelope
$[\tau_{\min}, \tau_{\max}] = [0.5,\, 3.0]$. By \Cref{eq:cps},

\[
S_{\mathrm{cool}} = \clamp{75 \times 0.95} = \mathbf{71.25}.
\]

The temperature range follows from \Cref{eq:dtmin,eq:dtmax}, and the clipping
rule binds on the lower bound:

\begin{align*}
\dT_{\min} &= \clip{0.5 \times 0.95}{0.5}{3.0} = \clip{0.475}{0.5}{3.0}
            = \mathbf{0.50}\ \si{\degreeCelsius}, \\
\dT_{\max} &= \clip{3.0 \times 0.95}{0.5}{3.0} = \clip{2.85}{0.5}{3.0}
            = \mathbf{2.85}\ \si{\degreeCelsius}.
\end{align*}

The raw product $0.5 \times 0.95 = 0.475$ falls below the literature floor and
is clipped up to it; the upper bound is unaffected because $A < 1$. This is the
point marked on \Cref{fig:clipping}.

The midpoint is $\overline{\dT} = (0.50 + 2.85)/2 = \qty{1.675}{\degreeCelsius}$,
which exceeds \qty{1.5}{\degreeCelsius} and therefore falls in the \emph{high}
heat-index improvement bucket. Shade potential, by \Cref{eq:shade}, is
$\clamp{1200/6000 \times 100} = \SI{20.0}{\percent}$.

\subsection{Layer 3 --- Impact and feasibility}

\paragraph{Energy.} All three preconditions of \Cref{sec:energy} hold: nearby
building cooling demand is confirmed relevant, an annual demand of
\qty{450000}{\kilo\watt\hour} is supplied, and street tree planting is flagged
as capable of affecting adjacent building loads. By \Cref{eq:energy-range},

\begin{align*}
E_{\min} &= 450\,000 \times 0.50 \times 0.02 = \qty{4500}{\kilo\watt\hour\per\year}, \\
E_{\max} &= 450\,000 \times 2.85 \times 0.04 = \qty{51300}{\kilo\watt\hour\per\year}.
\end{align*}

The upper bound is roughly eleven times the lower. That width is not a defect
of the arithmetic; it is the product of a sixfold temperature envelope and a
twofold sensitivity range, and it is the honest consequence of what the
evidence supports. A user who needs a tighter figure needs a building energy
model.

\paragraph{Greenhouse gases.} Version \methver{} ships no country emission
factors (\Cref{sec:ghg}). \emph{For this illustration only}, assume the
deploying authority has supplied a locally verified grid emission factor of
\qty{0.30}{\kilo\gram\per\kilo\watt\hour}. This value is not a Criterra
estimate, is not shipped with the tool, and carries no citation; it is a
placeholder input chosen to demonstrate \Cref{eq:ghg}. On that assumption,

\[
G = \qtyrange{1350}{15390}{\kilo\gram} \text{ CO\textsubscript{2}e per year,
or about \qtyrange{1.4}{15.4}{\tonne} CO\textsubscript{2}e per year.}
\]

Under the shipped configuration with no factor supplied, this output would read
\emph{not calculated}.

\paragraph{Costs.} No capital cost and no energy price were supplied.
Consistent with \Cref{sec:costs}, capital cost, annual cost savings, simple
payback, and cost feasibility are all reported as \emph{not estimated}. Cost
feasibility is therefore excluded from the final aggregation and its weight
redistributed.

\paragraph{Co-benefits.} Street tree planting carries library default levels of
\emph{medium} biodiversity, \emph{medium} stormwater, \emph{high} public
health, and \emph{high} urban quality. As noted in
\Cref{sec:composite-subscores}, the library supplies no default for social
inclusion, so it falls to the neutral \emph{unknown} value of 50 and is
itemised as an applied assumption. By \Cref{eq:cobenefit},

\begin{align*}
S_{\mathrm{cobenefit}}
&= 0.25(50) + 0.25(50) + 0.20(75) + 0.15(50) + 0.15(75) \\
&= 12.50 + 12.50 + 15.00 + 7.50 + 11.25 = \mathbf{58.75}.
\end{align*}

\paragraph{NbS Suitability.} Applying the rules of
\Cref{tab:suitability-rules}: the site offers sixty times the
\qty{100}{\metre\squared} minimum ($S_{\mathrm{space}} = 100$); soil
availability \emph{limited} exactly meets the typology's \emph{limited}
requirement ($S_{\mathrm{soil}} = 75$); irrigation \emph{occasional} exactly
meets the \emph{occasional} requirement ($S_{\mathrm{water}} = 75$);
maintenance intensity and land use were not supplied, so both take the neutral
50, itemised. Then

\[
S_{\mathrm{suit}} = 0.30(100) + 0.25(75) + 0.20(75) + 0.15(50) + 0.10(50)
                   = \mathbf{76.25},
\]

close to --- and slightly below --- the 78 assumed in the previous version of
this example, because the neutral defaults for maintenance and land use hold
the score down where full inputs would not.

\subsection{Aggregation}

Cost feasibility is unavailable, so \Cref{eq:aggregation-redistributed} applies
over the remaining five components, whose nominal weights sum to
$\sum_{i \in \mathcal{A}} w_i = 0.90$.

\begin{table}[H]
\centering
\caption[Worked example aggregation]{Aggregation of the illustrative
assessment. Cost feasibility is excluded and its \num{0.10} weight
redistributed proportionally across the remaining five components.}
\label{tab:we-aggregation}
\small
\begin{tabular}{@{}lrrrr@{}}
\toprule
\textbf{Component} & \textbf{Score} & \textbf{Nominal $w_i$} &
\textbf{Renormalised $w_i'$} & \textbf{Contribution} \\
\midrule
Heat Priority Index  & 71.75 & 0.25 & 0.2778 & 19.93 \\
Cooling Potential    & 71.25 & 0.25 & 0.2778 & 19.79 \\
NbS Suitability      & 76.25 & 0.15 & 0.1667 & 12.71 \\
Vulnerability        & 80.00 & 0.15 & 0.1667 & 13.33 \\
Co-benefits          & 58.75 & 0.10 & 0.1111 &  6.53 \\
Cost Feasibility     & \multicolumn{3}{l}{\emph{not estimated --- excluded}} & --- \\
\midrule
\textbf{NbS Cooling Opportunity Score} & & & & \textbf{72.29} \\
\bottomrule
\end{tabular}
\end{table}

The NbS Cooling Opportunity Score is \textbf{72.3}, which falls in the
\emph{Strong opportunity} band (61--80) of \Cref{tab:hpi-bands}.

\subsection{Confidence and reported result}

Applying the branched confidence model of \Cref{sec:confidence} with the
field-to-block mapping enumerated at version \methver{}: the cooling block has
eight of its ten slots supplied (\SI{80}{\percent} $\rightarrow$
\emph{high}; street tree planting carries \emph{high} evidence confidence, so
no cap applies). The energy block has the relevance confirmation and the demand
but no energy price (two of three, \SI{66.7}{\percent} $\rightarrow$
\emph{medium}). The economic block has none of its four fields
(\emph{low}). The equity block has three of its six fields
(\SI{50}{\percent} $\rightarrow$ \emph{medium}). The overall rating is the
lower median of \{high, medium, low, medium\}: \emph{medium}.

\begin{table}[htbp]
\centering
\caption[Worked example results]{The illustrative assessment as it would be
reported, with confidence ratings and applied assumptions. All temperature
values are daytime, pedestrian-level air temperature reductions.}
\label{tab:we-results}
\small
\begin{tabularx}{\textwidth}{@{}l>{\raggedright\arraybackslash}Xl@{}}
\toprule
\textbf{Output} & \textbf{Value} & \textbf{Confidence} \\
\midrule
Heat Priority Index        & 71.8 --- \emph{High} & High \\
Cooling Potential Score    & 71.3 & High \\
Temperature reduction      & \qtyrange{0.50}{2.85}{\degreeCelsius} (daytime, pedestrian-level air) & High \\
Heat-index improvement     & \emph{High} (midpoint \qty{1.675}{\degreeCelsius}) & High \\
Shade potential            & \SI{20.0}{\percent} of site & High \\
Energy savings             & \qtyrange{4500}{51300}{\kilo\watt\hour} per year & Medium \\
GHG avoided                & \qtyrange{1.4}{15.4}{\tonne} CO\textsubscript{2}e per year\textsuperscript{b} & Medium \\
Capital cost               & \emph{Not estimated --- no cost data supplied} & Low \\
Simple payback             & \emph{Not estimated} & Low \\
Cost feasibility           & \emph{Not estimated --- excluded from aggregation} & Low \\
Co-benefit Score           & 58.8 & Medium \\
\textbf{Opportunity Score} & \textbf{72.3 --- \emph{Strong opportunity}} & Medium \\
\bottomrule
\end{tabularx}
\end{table}

\noindent\footnotesize\textsuperscript{b}\,On an assumed locally supplied grid
emission factor of \qty{0.30}{\kilo\gram\per\kilo\watt\hour}. The shipped
configuration would report this as \emph{not calculated}.\normalsize

\medskip

\noindent\textbf{Assumptions the engine itemises for this run:} maintenance
intensity and land use defaulted to the neutral 50 (suitability); the grid
emission factor supplied by the user rather than from configuration; energy
price absent, so cost savings not computed; the four cited co-benefit library
defaults applied; social inclusion defaulted to the neutral 50 (no library
default); public accessibility, safety concern, and community participation
defaulted to the neutral 50 (equity); cost feasibility not estimated, excluded
from the aggregation with its weight redistributed.

\subsection{How much the result depends on the assumed values}
\label{sec:we-assumptions}

At the previous methodology version the NbS Suitability Score was assumed
rather than derived, and this subsection quantified the assumption's reach:
across the full 0--100 range of the sub-score the Opportunity Score spans
\numrange{67.9}{76.3} and never leaves the \emph{Strong opportunity} band. The
sub-score is now derived (\Cref{tab:suitability-rules}), and its derived value
of 76.25 lands within two points of the earlier assumption --- the conclusion
is unchanged. Category-migration behaviour across the whole scenario set is
quantified in \Cref{sec:sensitivity-results}.

The one remaining assumed input, the grid emission factor, affects only the
greenhouse-gas output, which scales linearly with it and enters no score.

\subsection{What this example demonstrates}

Four properties of the methodology are visible in the arithmetic above, and
each corresponds to a design commitment from \Cref{sec:commitments}.

\textbf{The clipping rule binds in practice.} The raw lower bound
$0.5 \times 0.95 = 0.475$ was clipped up to the literature floor. On a better
site the rule would bind at the ceiling instead, and the reported range would
not exceed \qty{3.0}{\degreeCelsius} regardless of how favourable the
conditions were.

\textbf{Refusal is a first-class output.} Three of the twelve reported
quantities are \emph{not estimated}, and a fourth --- annual cost savings ---
is not reported at all for want of an energy price. The tool did not substitute
a plausible
capital cost, and consequently did not produce a payback figure --- the single
number most likely to be extracted from a report and quoted without
qualification.

\textbf{Exclusion preserves ordering.} Because cost feasibility was excluded
rather than set to a neutral 50, this assessment remains comparable to another
option for the same site that also lacks cost data, without either being
rewarded for the absence.

\textbf{The ranges are wide, and stay wide.} The reported temperature range
spans nearly a factor of six and the energy range a factor of eleven. Nothing
in the methodology narrows them for presentational comfort.
